\section{Auxiliary Lemmas}\label{sec:auxiliary}

We use a regular supergraph to count the ways to extend a map by one
vertex. We refine the regularization lemma of Davey et
al.~\cite{DaveyEtAlStrong2026} for simple graphs by extending it to
loopless multigraphs. The construction preserves the multiplicities
of the original edges and gives every new edge multiplicity one.

\begin{lemma}\label{lem:regularization}
Let $D\geq1$ be an integer and let $G$ be a finite loopless multigraph
with $\Delta(G)\leq D$. There is a finite $D$-regular loopless multigraph
$\widehat G$ containing $G$ as an induced subgraph such that every
edge of $\widehat G$ outside $G$ has multiplicity one. In particular, $\widehat G$ is simple
if $G$ is simple.
\end{lemma}

\begin{proof}
We must add $D-d_G(v)$ incidences at each vertex $v$. Their total
is $D|V(G)|-2|E(G)|$. If it is odd, then $D$ is odd, so adjoining one
isolated vertex makes the total even. Form a list containing $D-d(v)$
occurrences of each vertex $v$ of the resulting graph, where $d(v)$ is
its current degree. Partition the list into pairs. For a pair $(u,v)$,
take a new copy of $K_{D+1}$, delete an edge $xy$, and add the edges
$ux$ and $vy$. Use a different copy for each pair. This adds two
distinct edges even when $u=v$. Both $x$ and $y$ regain degree $D$,
and every occurrence in the list adds one incidence at the
corresponding vertex. Thus all degrees are $D$. Each added edge has
multiplicity one and at least one new endpoint, so the original
vertices induce $G$.
\end{proof}

A strong clique in $G$ is also a strong clique in $\widehat G$, since
every adjacency between endpoints of its edges is preserved. In the
counting arguments we mark only the edges of the chosen clique, so
all edges added in the construction are unmarked.

Vector inequalities are coordinatewise. For symmetric matrices $X$
and $Y$ of the same size, $\tr(XY)=\sum_{i,j}X_{ij}Y_{ij}$.

In the simple-graph calculation we will express the size of a strong
clique, divided by the square of the maximum degree, as a linear
function of the pattern counts. The calculation uses a system with $N=23{,}789$ coordinates, $3{,}516$ linear equations and
$57$ positive semidefinite matrix constraints. Write $Aq=b$ for the
linear equations and $o^{\mathsf T}q$ for the objective, where
$A\in\mathbb Q^{3516\times N}$, $b\in\mathbb Q^{3516}$ and
$o\in\mathbb Q^N$. For each $\ell\in[57]$, let
$K_\ell,D_{\ell1},\ldots,D_{\ell N}$ be symmetric rational matrices of
the same size. The $\ell$th matrix constraint is $M_\ell(q)\succeq0$,
where
\begin{equation}\label{eq:intro-affine-gram}
 M_\ell(q)=K_\ell+\sum_{j=1}^N q_jD_{\ell j}.
\end{equation}
The matrices and vectors are specified in~\cite[Section~A]{FGSTSupplement}.
For a symmetric matrix
$Q$ of the same size as $K_\ell$, define
$D_\ell^*(Q)\in\mathbb R^N$ by
\begin{equation}\label{eq:intro-trace-adjoint}
 (D_\ell^*(Q))_j=\tr(QD_{\ell j})\qquad(j\in[N]).
\end{equation}

\begin{lemma}\label{prop:certificate-principle}
Let $o\in\mathbb Q^N$ be an objective vector and
$\lambda\in\mathbb Q^{3516}$. For each $\ell\in[57]$, let $Q_\ell$
be a matrix of the same size as $K_\ell$ expressible as a finite
nonnegative rational linear combination of outer products
$ww^{\mathsf T}$ of rational vectors. Set
\begin{align}
 &r=A^{\mathsf T}\lambda-\sum_{\ell=1}^{57}D_\ell^*(Q_\ell)-o,
 \label{eq:residual}\\
 &C=b^{\mathsf T}\lambda+\sum_{\ell=1}^{57}\tr(Q_\ell K_\ell).
 \label{eq:dual-value}
\end{align}
We call $r$ the residual vector. If $r\geq0$, then every $q\in\mathbb R^N$ with $q\geq0$, $Aq=b$ and
$M_\ell(q)\succeq0$ for all $\ell\in[57]$ satisfies $o^{\mathsf T}q\leq C$.
If $q\in[0,1]^N$ satisfies the same equations and matrix constraints,
then $o^{\mathsf T}q\leq C-\sum_{j:r_j<0}r_j$, for any residual $r$.
\end{lemma}

\begin{proof}
By \eqref{eq:intro-affine-gram} and
\eqref{eq:intro-trace-adjoint}, each trace
$\tr(Q_\ell M_\ell(q))$ equals
$\tr(Q_\ell K_\ell)+D_\ell^*(Q_\ell)^{\mathsf T}q$.
Adding the scalar product of $r$ in \eqref{eq:residual} with $q$
cancels the terms involving $D_\ell^*(Q_\ell)$. Since $Aq=b$,
we have $\lambda^{\mathsf T}Aq=b^{\mathsf T}\lambda$, so
the definition of $C$ in \eqref{eq:dual-value} gives
\begin{equation}\label{eq:certificate-identity}
 C-o^{\mathsf T}q
 =\sum_{\ell=1}^{57}\tr\bigl(Q_\ell M_\ell(q)\bigr)
  +\sum_{j=1}^N r_jq_j.
\end{equation}
Write $Q_\ell=\sum_a\alpha_{\ell a}w_{\ell a}w_{\ell a}^{\mathsf T}$,
where each $\alpha_{\ell a}$ is nonnegative. Then
\[
 \tr\bigl(Q_\ell M_\ell(q)\bigr)=
 \sum_a\alpha_{\ell a}w_{\ell a}^{\mathsf T}M_\ell(q)w_{\ell a}\geq0
\]
because $M_\ell(q)$ is positive semidefinite. If $r\geq0$ and $q\geq0$,
the second sum in \eqref{eq:certificate-identity} is nonnegative.
For $0\leq q_j\leq1$, we have $r_jq_j\geq0$ if $r_j\geq0$ and
$r_jq_j\geq r_j$ if $r_j<0$. Thus the second sum is at least
$\sum_{j:r_j<0}r_j$ in that case.
\end{proof}
